\needspace{20\baselineskip}
\section{mmap：只把 buffer 給 user space}
\label{sec:mmap}

\begin{oneline}
driver 的 \code{mmap} 只把 DMA buffer 那幾頁放進 app 的頁表，app 不必複製資料（zero-copy）就能讀寫 buffer，而暫存器從來不經由 driver 映射出去。
\end{oneline}

\begin{before}
v1.0.0 的 app 自己做兩次映射：暫存器（\code{/dev/uio4} 的 64 KB，第 1 章「原本怎麼做」的 \code{fuse_hal_init()}）和整塊實體記憶體（\code{/dev/mem} 從 \code{0x6000_0000} 起 256 MB，第 3 章「原本怎麼做」的 \code{fuse_arena_init()}）。板上 app 執行時，它的位址空間最後兩行就是這兩個映射；app 的任何一個錯誤指標，都可能寫到暫存器或那 256 MB 的任何地方：

\boardfile{A4_uio_state.txt}{5-12}
\end{before}

\begin{tech}
\begin{itemize}
\item \code{file_operations.mmap}：app 呼叫 \code{mmap()} 時，kernel 先建立一段位址範圍，再呼叫 driver 的這個函式，由 driver 決定這段範圍對應到哪些實體頁、用什麼屬性。
\item \code{struct vm_area_struct}（VMA）：描述這段範圍，包括起點 \code{vm_start}、終點 \code{vm_end}，以及 app 要求的 offset \code{vm_pgoff}（以頁為單位）。
\item \code{dma_mmap_coherent()}：把 \code{dma_alloc_coherent} 配的 buffer 映射到 user space，屬性和 kernel 端一樣是 non-cacheable。
\end{itemize}
\end{tech}

\begin{gain}
\begin{itemize}
\item app 的位址空間裡只剩一段 8 MB 的 \code{/dev/twn_fuse}，沒有暫存器，也沒有 \code{/dev/mem}。
\item 要求超過 buffer 大小，或 offset 不是 0，都會被拒絕（\code{-EINVAL}）。
\item 仍然是 zero-copy：輸入影像與 heatmap 不必經過 \code{read()}／\code{write()} 複製；CPU 端是 Normal non-cacheable，寫入可以合併（第 3 章）。
\end{itemize}
\end{gain}

\subsection{實際例子}

\textbf{driver 端：只映射 buffer。}先檢查 offset 與長度，再交給 DMA API。中間的 \code{map_device} 分支只給第 3 章的 A/B 實驗用：它改用 \code{remap_pfn_range()} 加 \code{pgprot_noncached()}，得到和 \code{/dev/mem} 一樣的 Device 屬性；有 SMMU 時 buffer 沒有固定的實體位址可以這樣映射，所以拒絕：

\snip{drv_mmap}

\code{file_to_td(f)} 從 \code{struct file} 找回 driver 的 \code{struct twn_dev}（第 5 章）。

\textbf{app 端：問大小、映射、回報「基底位址 0」。}

\snip{app_arena_init}

\code{*phys = 0}：回報基底位址 0，app 用「基底加 offset」算出的位址就變成 offset，正好是 ioctl 要的東西（為什麼這樣 app 就不用改，見第 9 章）。

\textbf{板上：app 的位址空間只剩 buffer。}pipeline app 執行中，看它的 \code{/proc/<pid>/maps} 與開著的檔案：

\boardfile{H10_driver_state.log}{56-58}

\code{0xffff95afd000 - 0xffff952fd000 = 0x800000}，正好 8 MB。和前面 v1.0.0 的輸出相比，\code{/dev/mem} 與 \code{/dev/uio4} 都不見了。越界的要求被拒絕：

\boardfile{H4_board_check.log}{10-12}

\subsection{原理}

\textbf{\code{mmap} 做了什麼。}app 呼叫 \code{mmap(NULL, len, ..., fd, offset)} 時，kernel 在 app 的位址空間找一段空的範圍，建立 VMA，記下長度與 offset（\code{vm_pgoff = offset / 4096}），再呼叫 driver 的 \code{.mmap} 填頁表。driver 回傳錯誤碼，整個 \code{mmap()} 就失敗。

\textbf{\code{dma_mmap_coherent} 的內部。}沒有 SMMU 時走 \code{kernel/dma/direct.c} 的 \code{dma_direct_mmap}：用 \code{dma_pgprot} 把屬性改成 non-cacheable（同一塊記憶體若被 kernel 和 user 用不同的 cache 屬性映射，CPU 可能看到不一致的資料），檢查範圍不超出 buffer，再用 \code{remap_pfn_range} 把這段連續的實體頁一次填進頁表。發佈映像打開了 SMMU，buffer 的實體頁不連續（第 3 章），改走 \code{drivers/iommu/dma-iommu.c} 的 \code{iommu_dma_mmap}，逐頁映射。對 driver 來說都是同一個呼叫。

\textbf{app 為什麼碰不到暫存器。}app 的位址空間裡有什麼，完全由 kernel 建的頁表決定。CSR 只在 kernel 裡做過 \code{ioremap}，driver 的 \code{.mmap} 只映射 buffer，所以 app 的頁表裡沒有任何一項指向 \code{0xA000_0000}。剩下的旁路是 root 自己開 \code{/dev/mem}：發佈映像用 \code{STRICT_DEVMEM} 與 \code{IO_STRICT_DEVMEM} 擋住 System RAM 與 driver 佔用中的暫存器，但有邊界（第 10 章）。

\textbf{映射的生命週期。}VMA 持有 \code{struct file} 的參考。就算 app 關掉 fd，只要映射還在，檔案就不算關閉，module 也就不能卸載（第 7 章的板上測試）。這保證 buffer 不會在 app 還在用的時候被釋放。

\subsection{面試會問}

\begin{qa}{app 為什麼只看得到 buffer，看不到暫存器？為什麼用 mmap 不用 read／write？}
app 的頁表完全由 kernel 決定。我的 driver 的 \code{.mmap} 只把 \code{dma_alloc_coherent} 配的 buffer 用 \code{dma_mmap_coherent} 映射給 app，並檢查長度不超過 buffer、offset 必須是 0；暫存器只在 kernel 裡 \code{ioremap}，從來沒有映射出去。板上 \code{/proc/<pid>/maps} 可以看到：UIO 版本有 256 MB 的 \code{/dev/mem} 和 64 KB 的 \code{/dev/uio4}，driver 版本只剩 8 MB 的 \code{/dev/twn_fuse}。資料走 mmap 是因為每幀幾百 KB，用 \code{read}／\code{write} 要經 \code{copy_from_user}／\code{copy_to_user} 多複製兩次；mmap 讓 app 和 PL 看同一塊記憶體。
\end{qa}

\subsection{關鍵函式速查}

\begin{center}
\small
\begin{tabularx}{\textwidth}{l l X}
\toprule
名稱 & 位置 & 做什麼 \\
\midrule
\code{twn_mmap()} & \code{twn_fuse.c}:\lineof{drv_mmap} & 檢查 offset 與長度，映射 buffer \\
\code{map_device} 的檢查 & \code{twn_fuse.c}:\lineof{drv_mmap_iommu} & 有 SMMU 時拒絕 Device 映射 \\
\code{fuse_arena_init()} & \code{pl_hal_fuse_drv.c}:\lineof{app_arena_init} & 問大小、\code{mmap}、回報基底 0 \\
\code{dma_mmap_coherent} & kernel & 映射 coherent buffer，屬性與 kernel 端一致 \\
\bottomrule
\end{tabularx}
\end{center}
